% TOP_PROBLEM: {"id": 3061, "title": "The permanent conjecture", "category_id": 2, "category": {"id": 2, "name": "combinatorics", "display_name": "Combinatorics", "description": "Counting problems, graph theory, discrete structures.", "slug": "combinatorics", "order_index": 2, "created_at": "2026-07-31T15:26:25.671Z"}, "status": "open", "problem_number": "OPG-151", "set_id": 12, "statusQualification": "The arbitrary base field, omitted in the short catalogue statement, is specified following its source discussion."}
% FIRST_POSTED: 2026-10-01
% ATTEMPT_STATUS: unresolved
% UnsolvedMath ID 3061; source catalogue code OPG-151
% !TeX program = lualatex
% GPT-6 Astra Ultra research attempt; independent partial work, 2026-10-01.
\documentclass[11pt,a4paper]{article}
\usepackage[margin=25mm]{geometry}
\usepackage{fontspec}
\setmainfont{lmroman10-regular.otf}[BoldFont=lmroman10-bold.otf,
 ItalicFont=lmroman10-italic.otf,BoldItalicFont=lmroman10-bolditalic.otf]
\usepackage{amsmath,amssymb,mathtools}
\usepackage{unicode-math}
\setmathfont{latinmodern-math.otf}
\AtBeginDocument{\let\setminus\smallsetminus}
\usepackage{xurl,hyperref}
\hypersetup{colorlinks=true,urlcolor=blue}
\setcounter{secnumdepth}{0}
\setlength{\parindent}{0pt}
\setlength{\parskip}{5pt}
\setlength{\emergencystretch}{3em}
\begin{document}
\section*{The permanent conjecture}\label{3061}
{\small UnsolvedMath ID 3061 · OPG-151 · Combinatorics · OpenGarden}
\subsection{Definitions and mathematical statement}
Let $\mathbb F$ be a field, let $n\ge1$, and let
$A=[a_1\ \cdots\ a_n]\in\mathbb F^{n\times n}$ have nonzero determinant.
Form the $n\times2n$ matrix $[A\ A]$ by placing two labelled copies of
each column side by side. An $n\times n$ submatrix in this problem retains
all rows and selects $n$ distinct column positions. Thus an original
column $a_j$ may appear zero, one or two times in the selected matrix,
but never more than twice.

For $B=(b_{ij})\in\mathbb F^{n\times n}$, its permanent is
\[
 \operatorname{per}(B)=\sum_{\sigma\in S_n}
                         \prod_{i=1}^n b_{i,\sigma(i)}.
\]
All terms have positive sign; their sum is evaluated in $\mathbb F$ and
can vanish through cancellation. Reordering selected columns does not
change the permanent.

The conjecture is that every such invertible $A$ admits a selection $B$
with $\operatorname{per}(B)\ne0$. The field is arbitrary, as specified
in the source discussion; no order, positivity, or characteristic-zero
hypothesis is intended. Equivalently, some multiplicity vector
$(m_1,\ldots,m_n)\in\{0,1,2\}^n$ with $\sum_jm_j=n$ gives a matrix
with nonzero permanent.

The task is only to establish this property for the duplicated matrix.
Replacing its second copy by a different invertible matrix is the
stronger generalization mentioned by the source, not an additional
requirement of this entry. Nonzero determinant of $A$ alone does not
identify the required column selection.
\subsection{Short English statement}
Does duplicating the columns of any invertible matrix always permit a square column selection with nonzero permanent?
\subsection{Research attempt: truncated polynomial rings in small characteristic}
\paragraph{Scope and literature check.}
GPT-6 Astra Ultra, 1 October 2026. The duplicated columns are labelled,
and a selected column type can occur at most twice. The original
discussion [S2] asks also about two different invertible matrices; that
stronger problem is not used here. The polynomial viewpoint is consistent
with Alon--Tarsi [S3]. The work below proves the duplicate-column claim
in characteristics two and three and explains the obstruction to a
straightforward extension.

\paragraph{1. Identify the relevant polynomial coefficients.}
Put
\[
 P_A(x_1,\ldots,x_n)=\prod_{i=1}^n\left(\sum_{j=1}^n a_{ij}x_j\right).
\]
For a multiplicity vector $m=(m_1,\ldots,m_n)$ with $\sum m_j=n$,
let $A_m$ repeat column $j$ exactly $m_j$ times. Expansion by which
column is assigned to each row shows
\[
 \operatorname{per}(A_m)=
 \left(\prod_j m_j!\right)[x_1^{m_1}\cdots x_n^{m_n}]P_A.
\]
Each assignment counted in the coefficient has exactly
$\prod_jm_j!$ labelled permutations of its equal-type columns. This
identity is an equality in the base field and does not silently divide
by a factorial.

\paragraph{2. Characteristic three.}
Let $R=\mathbb F[x_1,\ldots,x_n]/(x_1^3,\ldots,x_n^3)$.
Because $(\sum_j a_{ij}x_j)^3=\sum_j a_{ij}^3x_j^3$ in
characteristic three, the invertible substitution $x\mapsto Ax$
preserves the ideal. The inverse substitution does too, so it induces
an automorphism of $R$. The class of $x_1\cdots x_n$ is nonzero:
the residue monomials with every exponent in $\{0,1,2\}$ form a
vector-space basis of $R$. Its image, the class of $P_A$, is therefore
nonzero. Some coefficient with all $m_j\le2$ survives. All relevant
factorials are products of $1$ and $2$, hence are nonzero, and the
displayed coefficient identity gives an allowed nonzero permanent.
This proof works over every field of characteristic three, finite or not.

\paragraph{3. Characteristic two and a general comparison.}
In characteristic two the permanent equals the determinant, since
$-1=1$. Taking one copy of every column already works. More generally,
in characteristic $p$, the same truncated-ring proof with the ideal
$(x_1^p,\ldots,x_n^p)$ gives a nonzero permanent using at most $p-1$
copies of each column. Factorials of numbers less than $p$ are invertible.
This observation answers the duplicated question only for $p\le3$;
for larger $p$ it allows too many copies.

\paragraph{4. A check in every characteristic for size two.}
Write $A=\left(\begin{smallmatrix}a&b\\c&d\end{smallmatrix}\right)$.
If $ad+bc\ne0$, select the original columns. Otherwise invertibility
implies $ad-bc=-2bc\ne0$. Thus $2\ne0$, and $a,c\ne0$ because
$ad=-bc\ne0$. Two copies of the first column have permanent
$2ac\ne0$. This proves all $2\times2$ cases, without relying on
an order or positivity on the field.

\paragraph{Verification and the exact failed extension.}
The inverse-substitution step is essential; merely mapping the ideal
into itself would not by itself prove nonvanishing of an image. In
characteristic zero or $p>3$, an invertible change such as
$x_1\mapsto x_1+x_2$ need not preserve the cubic ideal:
$(x_1+x_2)^3$ has mixed terms $3x_1^2x_2+3x_1x_2^2$.
They survive in the quotient. Accordingly the automorphism argument
does not exist there. Nonzero $P_A$ in the ordinary polynomial ring
does not guarantee a monomial with every exponent at most two.
That bounded-exponent nonvanishing is the first remaining gap.

\subsection{Final status}
UNRESOLVED over arbitrary fields. The duplicate-column statement is
proved here in characteristics two and three and for $n\le2$; the
generalized $(p-1)$-copy result is explicitly weaker for larger primes.
These are restricted deductions, not a claim of a new universal theorem.

\subsection{Sources}
\begin{itemize}
\item[S1] ULAM.AI and the UnsolvedMath Contributors, supplied problems.json, record 3061 (OPG-151), OpenGarden collection. Catalogue identity and source transcription; CC BY 4.0.
\url{https://huggingface.co/datasets/ulamai/UnsolvedMath}.
\item[S2] Open Problem Garden, The permanent conjecture. Original problem and discussion; the mathematical formulation below is expanded from the supplied source record.
\url{http://www.openproblemgarden.org/op/the_permanent_conjecture}.
\item[S3] N. Alon and M. Tarsi, \emph{A nowhere-zero point in linear mappings}, Combinatorica 9 (1989), 393--395, original polynomial-method paper, checked 1 October 2026.
\url{https://web.math.princeton.edu/~nalon/PDFS/Publications/A%20nowhere-zero%20point%20in%20linear%20mappings.pdf}.
\end{itemize}
\end{document}
